\section{Stereographic Projection}\label{stereographic-projection}

Now that we have an understanding of how lines and triangles behave on a
sphere

As of now we have studied two-dimensional spherical geometry embedded on
the surface of a three-dimensional object, the sphere. This is obviously
very natural, yet it begs the question whether we could somehow flatten
the surface of the sphere.

We already use projections for our figures, where the sphere is
orthographically or perspectively projected onto the 2d image. Which our
brains are very used to seeing, we are able to reconstruct a 3d
interpretation in our minds eye.

The questions is whether there are other projections that might prove
more useful in comparing spherical geometry to euclidean and hyperbolic
geometry. We want our projection to have useful properties.

\begin{itemize}
\tightlist
\item
  The projection should be bijective in order for us to have one to one
  corespondace between points in the projection and points in the
  embedding.
\item
  it should preserve ``useful'' geometric properties

  \begin{itemize}
  \tightlist
  \item
    conformal (angle preserving)
  \item
    geodesics should map into ''intuitive'' counterpart
  \end{itemize}
\item
  it should come with a metric
\end{itemize}

\begin{center}\rule{0.5\linewidth}{0.5pt}\end{center}

Let \(\Phi: S \subset \mathbb{R}^3 \to \mathbb{R}^2\) be the projection
from the sphere to the plane and
\(\Phi^{-1}: \mathbb{R}^2 \to S \subset \mathbb{R}^3\) the unprojection
that from the plane to the spehere. Let \((x,y,z) \in S\) denote points
on the sphere and \((X,Y) \in \mathbb{R}^2\) points on the plane

\[\Phi(x,y,z) = \left(\frac{x}{1-z}, \frac{y}{1-z}\right) = (X, Y)\]

\[\Phi^{-1}(X, Y) = \left(\frac{2X}{X^2+Y^2+1}, \frac{2Y}{X^2+Y^2+1}, \frac{X^2 + Y^2 - 1}{X^2+Y^2+1}\right)= (x,y,z)\]

The stereographic projection maps great circles to

\begin{itemize}
\tightlist
\item
  straight lines through the origin, if the great circle intersects the
  north pole (projection point)
\item
  or circles otherwise
\end{itemize}

Proof. Let a great circle be given by its unit normal vector
\(\hat n = (n_X, n_y, n_z)\). Points on the great circle are given by

\[xn_x + yn_y + zn_z = 0\] we substitute \(x,y,z\) by unprotecting the
corresponding point on the plane \((X,Y)\)

\[\begin{split}
\Leftrightarrow 
n_x\frac{2X}{X^2+Y^2+1} + n_y\frac{2Y}{X^2+Y^2+1} + n_z\frac{X^2 + Y^2 - 1}{X^2+Y^2+1}&=0 \\[1em]
\Leftrightarrow
2Xn_x + 2Yn_y + n_z(X^2 + Y^2 - 1) &=0
\end{split}\]

If \(n_z = 0\), in which case the great circle must intersect the north
and south pole, the equation degenerates into a linear equation, giving
us points on a straight line through the origin.

\[\Leftrightarrow 2Xn_x + 2Yn_y = 0\]

If \(n \not= 0\) we divide by \(n_z\)

\[\Leftrightarrow 2X\frac{n_x}{n_z}+2Y\frac{n_y}{n_z} + X^2 + Y^2 - 1 = 0\]

We substitute \(a = 2\frac{n_x}{n_z}\) and \(b = 2\frac{n_y}{n_z}\) for
readablity and rearange to

\[\Leftrightarrow (X^2 + aX) + (Y^2 + bY) = 1\]

finally, by completing the square we end up with the equation of a
circle

\[\Leftrightarrow \left(X + \frac{a}{2} \right)^2 + \left(Y + \frac{b}{2}\right)^2 = 1 + \left(\frac{a}{2}\right)^2 + \left(\frac{b}{2}\right)^2\]

We can determine the exact properties the projected great cirlces

In the first case, the line given by the normal vector
\(\hat n’ = \begin{pmatrix}n_x\\ n_y\end{pmatrix}\)

\[\begin{split}
2Xn_x + 2Yn_y &= 0\\
\Leftrightarrow Xn_x + Yn_y &= 0 \\
\Leftrightarrow \begin{pmatrix}X\\ Y\end{pmatrix} \cdot \begin{pmatrix}n_x \\ n_y\end{pmatrix} &= 0
\end{split}\] In the second case, the circle can be described by it's
center \(C = (c_x, x_y)\) and radius \(r\)

The radius is given by \[\begin{split}
r &= 1 + \left(\frac{a}{2}\right)^2 + \left(\frac{b}{2}\right)^2 
= 1 + \frac{n_x^2+ n_y^2}{n_z^2}\\[1em]
c_x &= -\frac{a}{2} = -\frac{n_x}{n_z} \\[1em]
c_y &= -\frac{b}{2} = -\frac{n_y}{n_z}
\end{split}\]

\begin{itemize}
\item
  The circle is smallest with radius \(r=1\) when
  \(\hat n = (0, 0, \pm 1)\), i.e.~when the great circle lies on the
  equator
\item
  The radius grows to infinity as
  \(r \overset{n_z \rightarrow 0}\longrightarrow \infty\)
\item[$\square$]
  show how the formula changes if we offset the sphere along the z-axis,
  while still projecting from it's north pole
\item[$\square$]
  no isometric mapping
\item[$\square$]
  to get a better understanding of hyperbolic models in the next
  chapter, we want to understand the projection S2 onto E2
\item[$\square$]
  Great circles map to circles or lines
\item[$\square$]
  Reflection on lines becomes inversion
\item[$\square$]
  Excursion into mapping of transformations on the sphere
\item[$\square$]
  Distortions through the projection
\end{itemize}

\subsection{Polar Coordinates}\label{polar-coordinates}

\begin{itemize}
\tightlist
\item[$\square$]
  r: radius
\item[$\square$]
  \(\theta\): inclination
\item[$\square$]
  \(\phi\): azimuth (polar angle)
\item[$\square$]
  conversation cartesian coordinates
\end{itemize}

\subsection{Distance Metric}\label{distance-metric}

The distance between two points on a sphere is given by the minor arc
length of a great circle.

\begin{itemize}
\tightlist
\item[$\square$]
  exercise generalise the formula for a sphere with radius \(R\)
\item[$\square$]
  how we measure distance on the sphere
\item[$\square$]
  how we measure distance in our model
\item[$\square$]
  area metric
\item[$\square$]
  central angle \(\Delta \sigma = \mathbf a \cdot \mathbf b\)
\end{itemize}

\subsection{Area of a circle}\label{area-of-a-circle}

Let \(P\) be a point on the sphere. We define a circle with radius \(r\)
to be the set of points that are distance \(r\) away from \(P\)

The central angle is given by \(\Delta \sigma = \frac{R}{r}\).
